\section{Introduction}

A user asks an agent to create a meeting with Alice. After reading an untrusted
document, the agent calls the calendar tool with both Alice and Eve as
participants. The call creates the requested event and also commits an
attacker-chosen invitation, a common indirect-prompt-injection failure in
tool-using agents~\cite{agentdojo,toolemu}. The authorization question is not
whether the agent may invoke \texttt{create\_event}. It is whether this
execution may create this event, invite Alice, invite Eve, and apply the chosen
visibility.

Authorization cannot enforce a distinction that its observation interface has
already erased. An ACL may permit Alice and reject Eve, but a monitor that sees
only the tool name cannot apply that rule. Provenance may identify where the
string ``Eve'' originated without identifying whether it controls an
invitation, a message, or an inert note. Explicit arguments expose more of the
request, yet defaults, aliases, list expansion, and pre-state can change the
committed effect without changing the call text. Downstream model reasoning is
subject to the same information boundary.

\begin{figure}[t]
\centering
\begin{tikzpicture}[
  font=\scriptsize,
  call/.style={draw=TEBlue, fill=TEBlueLight, align=left,
               text width=3.45cm, minimum height=9mm, inner sep=3pt},
  effects/.style={draw=TEDark!70, fill=black!3, align=left,
                  text width=3.45cm, minimum height=12mm, inner sep=3pt},
  view/.style={draw=TEDark!85, align=left, text width=3.45cm,
               minimum height=15mm, inner sep=3pt},
  coarse/.style={view, draw=TERed, fill=TERedLight, dashed},
  atoms/.style={view, draw=TETeal, fill=TETealLight, line width=.7pt},
  flow/.style={-{Latex[length=1.7mm]}, draw=TEDark!85, line width=.5pt},
  guide/.style={draw=TEDark!55, line width=.4pt},
  stage/.style={circle, draw=TEBlue, fill=TEBlueLight,
               font=\scriptsize\bfseries, inner sep=1.2pt,
               minimum size=4.1mm},
  badge/.style={circle, draw=TEDark!55, fill=white,
                minimum size=6.4mm, inner sep=0pt},
]

\node[stage] at (-4.00,0) {1};
\node[call] (ucall) at (-2.0,0)
  {\textbf{Requested call $u$}\\[-1pt]
   \texttt{create\_event([Alice])}\\[-1pt]
   \hfill $P(u)=\Allow$};
\node[call] (vcall) at (2.0,0)
  {\textbf{Expanded call $v$}\\[-1pt]
   \texttt{create\_event([Alice,Eve])}\\[-1pt]
   \hfill $P(v)=\Deny$};
\node[badge] (ucal) at ($(ucall.north east)+(-.20,.05)$) {};
\path (ucal.center) pic[scale=.72] {te-calendar};
\node[badge, draw=TEOrange, fill=TEOrangeLight] (vcal)
  at ($(vcall.north east)+(-.20,.05)$) {};
\path (vcal.center) pic[scale=.72] {te-calendar};

\node[stage] at (-4.00,-1.72) {2};
\node[effects] (ueffects) at (-2.0,-1.72)
  {$\mu(u)=\{\textsf{create}(e),$\\
   \hspace*{9mm}$\textsf{invite}(e,A)\}$};
\node[effects] (veffects) at (2.0,-1.72)
  {$\mu(v)=\mu(u)\uplus$\\
   \hspace*{9mm}$\{\mathbf{\textsf{invite}(e,E)}\}$};
\node[badge] (ulayers) at ($(ueffects.north east)+(-.20,.05)$) {};
\path (ulayers.center) pic[scale=.72] {te-layers};
\node[badge, draw=TEOrange, fill=TEOrangeLight] (vlayers)
  at ($(veffects.north east)+(-.20,.05)$) {};
\path (vlayers.center) pic[scale=.72] {te-layers};

\draw[flow] (ucall) -- (ueffects);
\draw[flow] (vcall) -- (veffects);

\node[stage] at (-4.00,-3.75) {3};
\node[coarse] (coarse) at (-2.0,-3.75)
  {\textbf{Coarse call view}\\[-1pt]
   $\rho_c(u)=\rho_c(v)=\texttt{create\_event}$\\[1pt]
   \textit{Required decision is lost.}};
\node[atoms] (atomview) at (2.0,-3.75)
  {\textbf{Effect-atom view}\\[-1pt]
   $\rho_a(v)=\rho_a(u)\uplus\{\textsf{invite}(e,E)\}$\\[1pt]
   \textit{The policy can separate $u$ and $v$.}};
\node[badge, draw=TERed, fill=TERedLight] (warn)
  at ($(coarse.north east)+(-.20,.05)$) {};
\path (warn.center) pic[scale=.72] {te-warning};
\node[badge, draw=TETeal, fill=TETealLight] (shield)
  at ($(atomview.north east)+(-.20,.05)$) {};
\path (shield.center) pic[scale=.72] {te-shield};

\coordinate (fork) at (0,-2.72);
\draw[guide] (ueffects.south) |- (fork);
\draw[guide] (veffects.south) |- (fork);
\draw[flow] (fork) -| (coarse.north);
\draw[flow] (fork) -| (atomview.north);

\node[anchor=north, font=\scriptsize\bfseries, text=TERed] at (-2.0,-4.62)
  {$\rho_c$ collides};
\node[anchor=north, font=\scriptsize\bfseries, text=TETeal] at (2.0,-4.62)
  {$\rho_a$ preserves the distinction};
\end{tikzpicture}
\caption{An executable authorization-separating counterexample.  Calls $u$ and
$v$ use the same tool, but source execution shows that $v$ adds an invitation
that the policy rejects.  Their collision falsifies the coarse view for every
domain containing this pair. The typed-effect interface preserves the
distinction and can therefore expose it to policy.}
\label{fig:representation-collision}
\end{figure}


Figure~\ref{fig:representation-collision} illustrates this boundary. Let
$\rho$ be the execution representation available to an authorizer. If
$\rho(u)=\rho(v)$ while an admissible policy must decide $u$ and $v$
differently, every deterministic consumer of $\rho$ returns the same decision
for both. Allowing admits one unauthorized execution; denying or abstaining
withholds one authorized execution. Complete mediation and least privilege
therefore require an adequate object of mediation. We call this representation
obligation \emph{tool-effect binding}.

Agent harnesses make tool-effect binding a first-class systems problem. They
register evolving third-party tools from schemas and natural-language
descriptions, ask an untrusted model to synthesize calls, and resolve important
semantics inside the execution layer. A call may expand one list into several
recipients, choose a default mode, overwrite an existing resource, or create a
different transition under another pre-state. Unlike a security-typed API, the
harness is rarely given a trusted effect object that already captures these
distinctions. A model can propose such an object, but the proposal supplies
neither semantic evidence nor authority.

This paper asks whether a harness can test a proposed authorization interface
against the implementation it will mediate. The key evidence is asymmetric.
One executable pair with different policy decisions and an identical
representation falsifies that interface wherever the pair is reachable. A
successful registration instead creates a scoped conformance record: it names
the implementation, interventions, policy family, and unresolved tests bound
to the frozen descriptor. This asymmetry makes falsification the primary
security result and conformance an auditable deployment contract.

Our registration procedure treats a candidate descriptor as an untrusted
hypothesis. It executes a valid call and controlled variants in copied
sandboxes, varying fields, omissions, list cardinality, defaults, provenance,
or pre-state. A source oracle records the committed transition. A separately
supplied policy-separation oracle determines whether the semantic difference
requires another authorization decision. When policy separates two executions
that the descriptor merges, the pair becomes an executable failure certificate
and forces refinement. Invalid and unresolved interventions remain in the
validation record.

We instantiate this methodology with typed effect atoms. Each atom represents
one policy-relevant effect occurrence and retains the resource, target,
operation, provenance, or qualifier needed by the selected policy family.
List-valued calls expand when their members require independent decisions. The
method itself is representation-neutral: an arguments-plus-trusted-state view
can also be sufficient. Typed effects place tool-specific semantic adaptation
in a registered descriptor; a state-aware request leaves that adaptation to a
downstream consumer. We evaluate both paths and measure where they place state,
code, normalization, and runtime cost.

The evaluation follows the security argument from implementation behavior to
authorization. Source replay first establishes compound and state-dependent
effects. Executable separating pairs then expose collisions in tool-name,
raw-call, and common-field interfaces. Post-freeze experiments run 11 tools
from three public MCP implementations and compare candidate interfaces against
a native-delta decision chain. A four-domain authorization benchmark holds
calls, state, authority, and decision logic fixed while changing only the
observation interface. Finally, copied-sandbox and AgentDojo studies exercise
frozen interfaces at audited pre-commit boundaries.

This paper makes three contributions:
\begin{itemize}[leftmargin=*]
  \item \textbf{Authorization observability.} We formalize the upper bound that
  an observation interface places on every downstream authorizer and derive the
  unavoidable safety--availability consequence of a policy-separating
  collision.
  \item \textbf{Executable interface falsification.} We combine source-executed
  effect witnesses, independently supplied policy separation, and
  representation collisions into concrete failure certificates and a
  counterexample-guided registration discipline for existing tools.
  \item \textbf{Validated interfaces and consumption.} We falsify coarse views
  on project-local and independently implemented tools, evaluate typed-effect
  and state-aware interfaces under shared ACL, capability, and delegation
  authority, quantify their semantic-placement tradeoffs, and exercise a
  frozen descriptor at pre-commit integration points.
\end{itemize}

The central outcome is a security discipline: an authorization interface is an
implementation-grounded contract, not a declaration to be trusted on sight.
Executable counterexamples expose missing distinctions before downstream policy
depends on them.
